\documentclass[11pt]{article}
\usepackage{amsmath,amssymb,amsthm,mathtools,hyperref}

\begin{document}
\section*{Comparative biological timing under explicit phase conventions}
\label{prob:302}
\label{rank:114}
\par\smallskip

\noindent\textbf{Domain.} Biological timing; common-factor inference. \quad \textbf{Retained status/source context.} Open comparison of timing mechanisms. A single universal biological clock is a hypothesis requiring tests, rather than an established mathematical open conjecture.

\paragraph{Definitions and assumptions.} For a declared organism and condition $c$, let $Y_j(t)$ measure $m$ processes, such as developmental stage, metabolic state and circadian phase. Periodic phases take values on the circle; irreversible developmental coordinates do not. Restrict maps $h_j$ to prespecified smooth finite-dimensional families and seek an increasing time coordinate $\tau_c(t)$ with $\tau_c(0)=0$, $\tau_c(t_*)=1$ at a common calibration landmark, and bounded derivative. The observation law, sampling frequencies and normalized metric $d_j$ on each output space are fixed. Independent perturbation classes and held-out conditions are declared before fitting.

\paragraph{Problem statement.} Let $m_{jc}(t)$ be the metric-space mean of $Y_j(t)$ defined in the precision criterion below. Determine whether one timing factor is adequate:
    \[\max_{c,j}\sup_{0\leq t\leq H}
      d_j\bigl(m_{jc}(t),h_j(\tau_c(t))\bigr)\leq\varepsilon,\]
    with the same $h_j$ across validation conditions. If not, identify the smallest permitted number of timing coordinates and give an experimentally discriminating failure of the one-factor model.

\paragraph{Scope of adequacy.} The displayed criterion concerns population mean trajectories under the declared observation law. Acceptance does not by itself identify a shared physiological controller, individual-level synchronization or a shared joint stochastic law. Rejection concerns the declared model class, including its output maps, calibration and mean conventions. A mechanistic interpretation of an experiment requires the separately declared interventions and causal assumptions; it does not replace either required deliverable when the one-factor model fails.

\paragraph{Additional binding assumptions and resolution criterion.} An ordinary expectation of an angular phase is coordinate-dependent or undefined. Use the declared metric-space mean $m_{jc}(t)$: an ordinary mean for Euclidean outputs and a unique Fr\'echet mean, or an explicitly fixed selection from the mean set, for circular outputs. Specify any exclusion of antipodal or nonunique regimes. Fixed finite families and gauge conventions are needed to distinguish a common timing factor from arbitrary simultaneous reparameterization of $\tau$ and $h_j$.
For a circular output space $S_j$, define its Fr\'echet mean set by
\[
 F_{jc}(t)=\operatorname*{arg\,min}_{z\in S_j}
 \mathbb{E}\!\left[d_j(Y_j(t),z)^2\mid c\right],
\]
where the expectation uses the declared observation law. Thus the declared metric determines the objective; a geodesic metric and a chordal metric must not be interchanged silently. Require the mean set to exist wherever the criterion is evaluated. If it is nonunique, the selection rule is fixed before fitting and applied across times and conditions, rather than selected to improve the fit. For Euclidean outputs the ordinary mean must exist as a finite vector. A discrete or other nonconvex output space is not automatically a Euclidean output space merely because its states have numerical labels; its readout space and applicable mean convention must be declared.

\paragraph{Declarations needed for a concrete comparison.} Before fitting, declare the permitted coordinate counts and the candidate class at each count, including coordinate domains, output-map domains and families, calibration, symmetries or gauge conventions, complexity restrictions, any nesting, and how failure of every permitted count is reported. The scalar calibration alone does not specify this hierarchy. Also declare which quantities are learned during training, which observations may determine $\tau_c$ in a held-out condition, when they are available, and which observations are reserved for evaluation. These declarations distinguish retrospective representation from transfer or prediction. Specify coordinate regularity, the meaning of increasing, derivative bounds with units and their dependence on conditions, and the calibration landmark and its chronological timing relative to $H$; verify compatibility of these restrictions. Endpoint normalization alone need not remove nonlinear reparameterization freedom. The validation condition domain must be declared: the displayed maximum is retained, and replacing it by a supremum over a nonattaining domain would be a different convention requiring an explicit choice.

\paragraph{Alternative formulations.} (i) Within developmental processes, species-dependent time rescaling can test shared biochemical tempo. (ii) Circadian phase coupled to a developmental clock is a two-factor model, allowing independent entrainment. (iii) A monotone age coordinate with stationary hazard dependence addresses ageing rather than periodic oscillations. Equivalence requires actual shared scaling and coupling assumptions, not a relabeling of each observed trajectory.

\paragraph{Obstructions and partial progress.} Segmentation-clock experiments establish mechanism-specific interspecies timing differences; September 2026 proteomic work adds evidence linking protein stability and developmental tempo. These advances do not demonstrate common control of circadian, developmental and ageing clocks. Unrestricted output maps would make one-dimensional trajectory fitting nearly vacuous. Multiple sampling rates, phase slips and independently resettable processes challenge a shared coordinate. Model complexity limits and perturbation transfer make the proposed unification empirically falsifiable.

\paragraph{Selective external check in this revision.} SC08 (13 October 2026; access level and exact support are recorded in the source-check appendix). These selected checks are not an exhaustive present-day status search.

\paragraph{Sources (supplied provenance; not all rechecked).} \begin{itemize}
\item Supplied Wikipedia background: unsolved biology list. \url{https://en.wikipedia.org/wiki/List_of_unsolved_problems_in_biology} (Retained for provenance; not used as scientific verification of current open status.)
\item Matsuda et al., Species-specific segmentation clock periods are due to differential biochemical reaction speeds (2020). \url{https://pubmed.ncbi.nlm.nih.gov/32943519/} (Primary experimental abstract checked; developmental timing only.)
\item Systematic differences in protein stability underlie species-specific developmental tempo (2026). \url{https://pubmed.ncbi.nlm.nih.gov/42594862/} (Indexed primary abstract checked; published 9 September 2026, not a universal timing theorem.)
\end{itemize}


\section*{Binding catalogue conventions}
This edition revises both supplied catalogues and preserves original problem numbers 1--330. It adds one globally unique label \verb|\label{rank:x}| per problem, where smaller $x$ means greater provisional difficulty. A rank label is an anchor at that problem's section: \verb|\ref{rank:x}| returns the section number, not the literal difficulty number. The visible \texttt{rank:x} text and the index give the difficulty number itself.

The entries are \emph{parameterized mathematical research schemes}, including empirical prediction, causal identification, historical reconstruction, theory and engineering. They are not 330 certified unsolved mathematical conjectures. A rigorous instance fixes the objects and restrictions below; changing those objects changes the question. Established identities, finite calculations, non-identification constructions and solved special cases are not being posed anew as open. Some residual tasks concern informative bounds, biological realization or external validation rather than a new theorem.

Every main definition and problem statement received a mathematical editorial review. Each entry now has an additional binding-conditions paragraph; direct corrections are separately recorded in the accompanying audit. The alternative formulations inherit the applicable conventions, but are different targets unless equivalence is proved. Where an old broad phrase conflicts with an explicit restriction below it, the explicit restriction controls.

\textbf{Literature boundary.} Thirteen selective primary-literature checks cover eighteen entries in this revision; their access levels and scope are recorded explicitly. The supplied bibliography and its earlier source-access descriptions are retained as provenance, not represented as newly verified. Neither this review nor a source's description proves the absence of later solutions. The date of this edition is 13 October 2026; there is no claim of exhaustive current-status certification for all 330 entries.

\subsection*{Binding conventions for a concrete instance}
\paragraph{Objects, quantifiers and non-vacuity.}
An instance declares measurable state, input, output and parameter spaces; a fixed candidate family; admissible initial laws, interventions and environments; a finite horizon or specified asymptotic regime; observation and sampling laws; and performance metrics, positive normalization scales and positive tolerances. These are chosen before fitting or testing, not adjusted after seeing failures. Shared parameters must genuinely be shared across the declared validation domain. Finite dimension alone does not restrict arbitrary real-number encodings: specify a finite library, bounded description length, or regularity/complexity bounds. Exclude constant or no-signal solutions through a fixed informative-output or baseline-improvement requirement where appropriate. A universally quantified result cannot be established by choosing one trivial favorable instance.

\paragraph{Probability laws and data.}
Every expectation or probability refers to the specified stochastic law, including initial uncertainty, process noise, observation noise and sampling as appropriate. A deterministic equation is not a probability model without a declared random input or initial law. ``Observed'' targets in population-level formulas mean population observables $g_j^*$, not noisy estimates asserted to be exact. Let $\widetilde g_j$ be independent validation estimates and suppose
\[
\mathbb{P}\{\lvert\widetilde g_j-g_j^*\rvert\le r_j\text{ for every }j\}\ge1-\alpha.
\]
Then an uncertainty-aware sufficient certification rule for positive target tolerances $\varepsilon_j$ is
\[
\max_j\frac{|\widehat g_j-\widetilde g_j|+r_j}{\varepsilon_j}\le1.
\]
On that simultaneous confidence event it implies every population error is at most its tolerance, by the triangle inequality. Vector/norm versions use corresponding valid simultaneous bounds. Training uncertainty must also be included when the claim averages over repeated training data. Evaluation on finitely many conditions does not establish a continuous-domain supremum without covering and regularity bounds.

\paragraph{Dynamics and well-posedness.}
For ODEs and SDEs give coefficient domains, initial laws and existence/uniqueness conditions in the intended solution sense, or keep the entire admitted solution set. For stochastic jump models specify nonnegative off-diagonal rates, diagonal normalization, rate orientation and non-explosion. For PDEs specify spatial domains, boundary/initial conditions, conservation or source terms, and a weak or classical solution class. Moving-domain concentration equations include dilution/transport terms. A stationary law is used only for a time-homogeneous augmented model under existence assumptions; a stationary approximation to a driven system must be justified. Absorbing extinction states are retained. Survival-conditioned, quasi-stationary and unconditional laws are distinct and must not be substituted for one another.

\paragraph{Identification and causal meaning.}
Two parameter values are observationally equivalent when they induce the same full distribution under all allowed observation/intervention designs, not merely the same likelihood value for one sample. A target is identified only if constant on the equivalence class. Gauge symmetries, label switching and reference conventions are quotiented or fixed. Finite-data compatible sets use a declared confidence, likelihood or discrepancy rule; an empty set signals model rejection. Bayesian credible sets are not automatically confidence sets. A prior or regularizer does not create data identification. Bounds intended to resolve a question must be informative at a predeclared width or decision threshold.

A causal contrast uses well-defined intervention versions and a stated causal model. Random assignment, or explicit consistency, overlap and exchangeability assumptions, is required for causal interpretation from observations. Interference and time-varying confounding are represented when present. Fixing a mediator by intervention differs from conditioning on its observed value. Marginal potential-outcome laws do not determine individual cross-world effects or an arbitrary joint coupling. Ethically admissible designs, rather than unconstrained interventions, define the experimental menu.

\paragraph{Metrics, units, optimization and limits.}
Specify the sampling or reference measure for densities and integrals, the log base, and dimensionless arguments of logarithms. Total variation is $\sup_A|P(A)-Q(A)|$; Wasserstein distances require the declared ground metric and appropriate moment conditions; KL divergence uses its stated absolute-continuity convention. All normed readouts are scaled consistently, and error probabilities use $0<\alpha<1$. A ratio has a positive denominator or an explicit undefined/zero convention. Thermodynamic entropy production is nonnegative in expectation under the relevant hypotheses; individual stochastic realizations need not be. An infimum is not an attained minimum without compactness/coercivity and continuity or another existence argument. Fix the order of large-population, long-time and vanishing-mutation limits. Stability claims declare perturbation space and symmetries.

\subsection*{How the provisional difficulty order was assigned}
Four editorial components are scored from 1 (lower burden) to 5 (higher burden): (a) formal or inferential difficulty, (b) access to evidence and identification burden, (c) mechanistic/computational integration, and (d) independent validation and transfer. Their unweighted total orders the entries. For selected entries in the highest bands, an explicit editorial priority breaks equal totals; all other ties are broken by original problem number. Equal totals are \emph{not} evidence of reliably different true difficulties. Scores are judgments, not measurements, completion-time estimates or consensus rankings.

The object ranked is an informative, nondegenerate resolution of the \emph{scoped} scientific target, with the stated budget and validation domain. A sharp useful impossibility result can qualify; an already known identity, vacuous bound, arbitrary favorable model or lookup table cannot. Rank 330 is merely the least difficult relative assessment in this catalogue, not a claim of an easy biological solution. Broad historical titles do not automatically outrank precisely formulated mechanistic tasks.


\section*{Selective source-check ledger}
\subsection*{SC08}\label{check:SC08}

Matsuda et al., Systematic differences in protein stability underlie species-specific developmental tempo (Developmental Cell, online 13 August 2026; 9 September 2026 issue), DOI 10.1016/j.devcel.2026.07.012.

\textbf{Applied here to problems:} 132, 302.

\textbf{Access level:} Primary-paper abstract and exact online/issue dates.

\textbf{Support and limit:} Protein-stability differences and perturbations support a mechanism in the studied human and mouse developmental systems, not one universal clock for all biological processes.

\url{https://pubmed.ncbi.nlm.nih.gov/42594862/}


\end{document}
